% !TeX program = pdflatex
% !BIB program = biber
% Modelo: TCC completo com a classe abntex2 (arquivo principal e um arquivo por capítulo)
% Compilar: latexmk -pdf tcc.tex   Testado com: abntex2, biblatex-abnt, acro 3, TeX Live 2024+
% Compile sempre este arquivo, nunca um capítulo. Estrutura da pasta:
%   tcc.tex            preâmbulo, dados do trabalho e ordem dos elementos (NBR 14724)
%   pretextuais/       folha de aprovação e resumos (incluídos com \input)
%   capitulos/         um arquivo por capítulo, apêndice e anexo (incluídos com \include)
%   referencias.bib    base de referências (biblatex-abnt)
% A variante com a classe abntex2-ufv está em ../tcc-abntex2-ufv/.
%%BEGIN:modeloclasse
\documentclass[12pt, a4paper, oneside, english, brazilian,
  sumario=tradicional]{abntex2}
%%END:modeloclasse

%%BEGIN:modelopacotes
% 1. PACOTES. Primeiro a codificação e a fonte; o biblatex perto do fim.
% O hyperref já vem com a abntex2: não o carregue de novo.
\usepackage[T1]{fontenc}
\usepackage{lmodern}
\usepackage{indentfirst}          % recua o primeiro parágrafo da seção
\usepackage{graphicx}             % figuras
\usepackage{booktabs}             % filetes das tabelas
\usepackage{tabularx}             % colunas X nos quadros
\usepackage{xcolor}
\usepackage{listings}             % código-fonte
\usepackage{acro}                 % siglas e abreviaturas
\usepackage{csquotes}             % exigido pelo biblatex
\usepackage[backend=biber, style=abnt, uniquename=init, giveninits,
            repeatfields]{biblatex}
\usepackage{ragged2e}             % \RaggedRight com hifenização
\usepackage{xurl}                 % quebra URLs longas
\addbibresource{referencias.bib}
\appto{\bibsetup}{\RaggedRight}   % referências à esquerda (NBR 6023)
\urlstyle{same}
\hypersetup{hidelinks}            % links sem moldura (para imprimir)
%%END:modelopacotes

%%BEGIN:modeloflutuantes
% 2. QUADRO E CÓDIGO. A abntex2 não define o quadro: o float e a lista
% vêm do memoir, e a linha da lista imita a da lista de figuras.
\newfloat{quadro}{loq}{Quadro}
\newlistof{listofquadros}{loq}{Lista de quadros}
\newlistentry{quadro}{loq}{0}
\renewcommand{\cftquadroname}{Quadro\space}
\renewcommand*{\cftquadroaftersnum}{\hfill\textendash\hfill}
\setlength{\cftquadronumwidth}{\cftfigurenumwidth}
\providecommand{\quadroautorefname}{Quadro}
% uma lista por tipo: "Lista de figuras" (a abntex2 a chama de ilustrações)
\AtBeginDocument{\renewcommand{\listfigurename}{Lista de figuras}}
% O código usa a legenda do próprio listings, numerada sem o capítulo.
\lstset{basicstyle=\ttfamily\small, columns=flexible, keepspaces=true,
  numbers=left, numberstyle=\tiny, frame=single, breaklines=true,
  numberbychapter=false,
  literate={á}{{\'a}}1 {é}{{\'e}}1 {í}{{\'i}}1 {ó}{{\'o}}1 {ú}{{\'u}}1
           {ã}{{\~a}}1 {õ}{{\~o}}1 {ç}{{\c c}}1 {ê}{{\^e}}1 {â}{{\^a}}1}
\renewcommand{\lstlistingname}{Código}
\providecommand{\lstlistingautorefname}{Código}
\newlistof{listofcodigos}{lol}{Lista de códigos}  % \listofcodigos*
%%END:modeloflutuantes

%%BEGIN:modelosiglas
% 3. SIGLAS (acro): declare aqui e use \ac{chave} no texto. A lista sai
% com o título no padrão dos elementos pré-textuais.
\NewAcroTemplate[heading]{abnt}{\pretextualchapter{\acrolistname}}
\acsetup{list/name = {Lista de siglas}, list/heading = abnt,
  list/template = tabular,
  templates/colspec = {@{}l@{\hspace{1.5em}}p{.8\linewidth}@{}}}
\DeclareAcronym{abnt}{short = ABNT,
  long = Associação Brasileira de Normas Técnicas}
\DeclareAcronym{api}{short = API,
  long = Interface de Programação de Aplicações,
  long-plural-form = Interfaces de Programação de Aplicações}
\DeclareAcronym{csv}{short = CSV,
  long = \textit{Comma-Separated Values}}
%%END:modelosiglas

%%BEGIN:modelodados
% 4. DADOS DO TRABALHO (capa, folha de rosto e folha de aprovação)
\titulo{Título do trabalho}
\autor{Nome do Autor}
\local{Cidade}
\instituicao{Nome da Instituição \par Nome do curso}
\data{\the\year}
\orientador{Prof. Dr. Nome do Orientador}
\tipotrabalho{Trabalho de Conclusão de Curso (Graduação)}
\preambulo{Trabalho apresentado à Nome da Instituição como parte das
exigências para a obtenção do título de Bacharel em Nome do Curso.}

% Para compilar só os capítulos em que está trabalhando, descomente:
% \includeonly{capitulos/metodologia}
%%END:modelodados

\begin{document}

%%BEGIN:modelopre
% ELEMENTOS PRÉ-TEXTUAIS, na ordem da NBR 14724
\pretextual
\imprimircapa
\imprimirfolhaderosto  % a linha seguinte fica em branco (ele lê um argumento)

% Folha de aprovação (incluída por tcc.tex). Depois da defesa, troque esta página pelo PDF
% assinado: \includepdf{folha-assinada.pdf} (pacote pdfpages, no preâmbulo).
\begin{folhadeaprovacao}
  \begin{center}
    {\ABNTEXchapterfont\large\imprimirautor}

    \vspace*{\fill}
    {\ABNTEXchapterfont\bfseries\Large\imprimirtitulo}
    \vspace*{\fill}

    \hspace{.45\textwidth}
    \begin{minipage}{.5\textwidth}
      \imprimirpreambulo
    \end{minipage}%
    \vspace*{\fill}
  \end{center}

  Trabalho aprovado. \imprimirlocal, dia de mês de ano.

  \assinatura{\textbf{\imprimirorientador} \\ Orientador}
  \assinatura{\textbf{Professor} \\ Convidado 1}
  \assinatura{\textbf{Professor} \\ Convidado 2}

  \begin{center}
    \vspace*{0.5cm}
    {\large\imprimirlocal}
    \par
    {\large\imprimirdata}
    \vspace*{1cm}
  \end{center}
\end{folhadeaprovacao}
   % \input: sem quebra de página
% Resumo e abstract (incluídos por tcc.tex). Parágrafo único, de 150 a 500 palavras; palavras-chave
% com ponto e vírgula, inicial minúscula e ponto final (NBR 6028:2021). Resumo
% sem siglas com \ac: ele é lido sozinho, e a primeira ocorrência por extenso ficaria aqui.
\begin{resumo}
Texto do resumo em português, em parágrafo único: objetivo, método,
resultados e conclusões do trabalho.

\textbf{Palavras-chave}: LaTeX; trabalho acadêmico; normas ABNT.
\end{resumo}

\begin{resumo}[Abstract]
\begin{otherlanguage*}{english}
Abstract text in English, in a single paragraph.

\textbf{Keywords}: LaTeX; academic writing; ABNT standards.
\end{otherlanguage*}
\end{resumo}


\listoffigures*        % o * deixa a lista fora do sumário
\cleardoublepage
\listofquadros*
\cleardoublepage
\listofcodigos*
\cleardoublepage
\listoftables*
\cleardoublepage
\printacronyms         % só as siglas usadas no texto
\cleardoublepage
\tableofcontents*
\cleardoublepage
%%END:modelopre

%%BEGIN:modelotexto
% ELEMENTOS TEXTUAIS: \textual antes do primeiro capítulo (numeração
% arábica e cabeçalhos). Cada \include começa em página nova.
\textual
% Capítulo incluído por tcc.tex com \include: sem preâmbulo e sem \begin{document}.
\chapter{Introdução}\label{cap:introducao}

A elaboração de trabalhos acadêmicos segue as normas da \ac{abnt}. Este
modelo mostra as peças de um trabalho completo: siglas, figura, quadro,
código, tabela, equação, citações e referências, cada capítulo em um
arquivo.

\section{Objetivos}\label{sec:objetivos}

O objetivo geral é desenvolver e avaliar um protótipo. Os objetivos
específicos são:
\begin{alineas}
  \item levantar os requisitos (\autoref{cap:metodologia});
  \item desenvolver o protótipo;
  \item avaliar o protótipo com usuários (\autoref{cap:resultados}).
\end{alineas}

\section{Organização do trabalho}

O \autoref{cap:referencial} apresenta os conceitos usados no trabalho; o
\autoref{cap:metodologia}, os métodos; o \autoref{cap:resultados}, os
resultados; e o \autoref{cap:conclusao}, as conclusões.
        % sem a extensão .tex
% Capítulo incluído por tcc.tex com \include: sem preâmbulo e sem \begin{document}.
\chapter{Referencial teórico}\label{cap:referencial}

Segundo \textcite{ex-livro}, a citação direta longa (mais de três linhas)
fica recuada:
\begin{citacao}
Texto citado literalmente, com mais de três linhas, que a classe abntex2
formata com recuo de 4\,cm, fonte menor e espaçamento simples, como orienta a
NBR 10520. A referência vem ao final da citação, entre parênteses
\cite[p.~35]{ex-livro}.
\end{citacao}

Uma \ac{api} define como dois sistemas trocam dados; as \acp{api} da
\autoref{fig:arquitetura} usam arquivos \ac{csv} \cite{ex-artigo}.

% Título acima, fonte e nota abaixo, na largura da imagem (minipage)
\begin{figure}[htbp]
  \centering
  \begin{minipage}{.6\linewidth}
    \caption{Arquitetura do protótipo}\label{fig:arquitetura}
    \includegraphics[width=\linewidth]{example-image-a}
    \fonte{Elaborado pelo próprio autor.}
  \end{minipage}
\end{figure}

% Capítulo incluído por tcc.tex com \include: sem preâmbulo e sem \begin{document}.
\chapter{Metodologia}\label{cap:metodologia}

Os requisitos estão no \autoref{qua:requisitos}, e o \autoref{lst:media}
mostra a função que calcula a média das avaliações.

% Quadro: conteúdo textual e bordas fechadas (float definido em tcc.tex)
\begin{quadro}[htbp]
  \caption{Requisitos funcionais do protótipo}\label{qua:requisitos}
  \renewcommand{\arraystretch}{1.25}
  \begin{tabularx}{\textwidth}{|c|X|c|}
    \hline
    \textbf{ID} & \textbf{Descrição} & \textbf{Prioridade} \\
    \hline
    RF01 & Cadastrar usuário com nome, \emph{e-mail} e senha. & Alta \\
    \hline
    RF02 & Exportar as avaliações em \acs{csv}. & Média \\
    \hline
  \end{tabularx}
  \fonte{Elaborado pelo próprio autor.}
\end{quadro}

% Código: a legenda vem do próprio listings (Lista de códigos em tcc.tex)
\begin{lstlisting}[language=Python, caption={Média das avaliações},
    label={lst:media}]
def media(notas):
    """Devolve a média; 0 para a lista vazia."""
    return sum(notas) / len(notas) if notas else 0.0
\end{lstlisting}
\fonte{Elaborado pelo próprio autor.}

% Capítulo incluído por tcc.tex com \include: sem preâmbulo e sem \begin{document}.
\chapter{Resultados e discussão}\label{cap:resultados}

A \autoref{tab:avaliacao} resume a avaliação, e a nota média $\bar{x}$ foi
calculada pela \autoref{eq:media}:
\begin{equation}\label{eq:media}
  \bar{x} = \frac{1}{n} \sum_{i=1}^{n} x_i
\end{equation}

% Tabela no padrão IBGE: título, fonte e nota na largura da tabela
\begin{table}[htbp]
\IBGEtab{\caption{Avaliação do protótipo por grupo de usuários}%
         \label{tab:avaliacao}}{%
  \begin{tabular}{lrr}
    \toprule
    \textbf{Grupo} & \textbf{Participantes} & \textbf{Nota média} \\
    \midrule
    Estudantes  & 18 & 8,4 \\
    Professores &  6 & 7,9 \\
    \bottomrule
  \end{tabular}%
}{\fonte{Elaborado pelo próprio autor.}%
  \nota{Dados fictícios, apenas para exemplo.}}
\end{table}

% Capítulo incluído por tcc.tex com \include: sem preâmbulo e sem \begin{document}.
\chapter{Conclusão}\label{cap:conclusao}
% Anexo: \autoref escreveria "Apêndice" (a abntex2 usa o mesmo contador);
% escreva a palavra e use \ref.

As conclusões respondem aos objetivos da \autoref{sec:objetivos}. O
questionário aplicado está no \autoref{ap:questionario}, e o regulamento do
curso, no Anexo~\ref{an:regulamento}.

%%END:modelotexto

%%BEGIN:modelopos
% ELEMENTOS PÓS-TEXTUAIS
\postextual
\printbibliography[title=Referências]

\begin{apendicesenv}                  % "APÊNDICE A – ..."
\partapendices
% Apêndice (texto do próprio autor), incluído dentro de apendicesenv em tcc.tex.
\chapter{Questionário aplicado}\label{ap:questionario}

Perguntas feitas aos participantes da avaliação.

\end{apendicesenv}

\begin{anexosenv}                     % "ANEXO A – ..."
\partanexos
% Anexo (documento de terceiros), incluído dentro de anexosenv em tcc.tex.
\chapter{Regulamento do curso}\label{an:regulamento}

Texto do regulamento, elaborado pela coordenação do curso.

\end{anexosenv}
%%END:modelopos

\end{document}
